\documentclass[11pt]{article}
\usepackage[margin=1in]{geometry}
\usepackage[T1]{fontenc}
\usepackage[utf8]{inputenc}
\usepackage{lmodern}
\usepackage{microtype}
\usepackage{array}
\usepackage{booktabs}
\usepackage{enumitem}
\usepackage{hyperref}
\usepackage{xcolor}
\hypersetup{
  colorlinks=true,
  linkcolor=black,
  urlcolor=blue!50!black,
  pdftitle={Coding Scheme Dossier}
}
\setlist[itemize]{leftmargin=1.6em}
\setlist[description]{style=nextline,leftmargin=3.5cm,labelsep=0.6em}
\newcommand{\field}[1]{\nolinkurl{#1}}
\newcommand{\filepath}[1]{\nolinkurl{#1}}
\newcommand{\schemeheader}[3]{
  \begin{center}
    {\LARGE \textbf{#1}}\\[0.5em]
    {\large #2}\\[0.4em]
    {\normalsize \textit{#3}}
  \end{center}
  \vspace{0.8em}
}



\begin{document}
\schemeheader
  {Scheme 6}
  {BPS Ontology and Semantic Loading Benchmark Scheme}
  {Ontology prompts for benchmark-relative semantic quantification}

\section*{Purpose}
This scheme supplies the ontology scaffold used to quantify semantic emphasis across biological, psychological, and social axes. It is not a manual adjudication form, but it is an operational text-classification framework because it standardizes the domain and subdomain prompts against which review records are embedded and compared.

\section*{Operational Status}
\begin{itemize}
\item Workflow position: semantic loading and synthesis stage after Stage 2 coding.
\item Operational mode: ontology-prompted embedding benchmark, with OpenRouter embeddings when available and TF-IDF fallback otherwise.
\item Canonical implementation basis: \filepath{src/bps_review/reporting/semantic_loading.py} and the archived ontology terms JSON.
\end{itemize}

\section*{Canonical Source Paths}
\begin{itemize}
\item \filepath{src/bps_review/reporting/semantic_loading.py}
\item \filepath{src/vector_db/semantic_loading/ontology/ontology_terms.json}
\item \filepath{src/vector_db/semantic_loading/analysis/domain_loading_summary.csv}
\item \filepath{src/vector_db/semantic_loading/analysis/subdomain_loading_summary.csv}
\end{itemize}

\section*{Record Text Used for Scoring}
Each review record is composed into a semantic analysis string with the following sections: title, abstract, and extracted objective text. The ontology scheme then compares that composed record text against the benchmark prompts below.

\section*{Prompt Templates}
\begin{description}
\item[\field{domain prompt}] \texttt{\{domain\} chronic pain ontology. \{joined subdomains\}.}
\item[\field{subdomain prompt}] \texttt{\{domain\} chronic pain subdomain. \{term\}.}
\end{description}

\section*{Domain Order}
The benchmark axes are always ordered as \texttt{biological}, \texttt{psychological}, and \texttt{social}.

\section*{Ontology Content}
\subsection*{Biological Subdomains}
\begin{itemize}
\item Central Sensitization and Neuroplasticity
\item Musculoskeletal and Structural Pathology
\item Pharmacological and Biomedical Treatment
\item Sleep Disruption and Circadian Dysregulation
\item Immune Inflammatory and Neuroinflammatory Processes
\item Genetic Epigenetic and Biological Vulnerability
\item Neuroimaging Brain Structure and Function
\item Nociceptive Transmission and Pain Pathways
\item Physical Function Mobility and Deconditioning
\item Metabolic Nutritional and Hormonal Factors
\end{itemize}

\subsection*{Psychological Subdomains}
\begin{itemize}
\item Catastrophizing and Negative Cognitive Appraisal
\item Fear Avoidance and Pain Related Fear
\item Depression Emotional Distress and Affect
\item Anxiety and Psychological Reactivity
\item Self Efficacy Control Beliefs and Perceived Mastery
\item Acceptance Psychological Flexibility and Mindfulness
\item Pain Coping Strategies and Adjustment
\item Attention Vigilance and Pain Processing
\item Illness Beliefs Pain Representations and Meaning
\item Cognitive Behavioral and Psychotherapeutic Approaches
\item Third Wave Therapies ACT and Contextual Approaches
\item Resilience Positive Psychology and Post Traumatic Growth
\item Identity Self Concept and Chronic Pain Biography
\item Trauma Adverse Childhood and Life Events
\item Personality Psychological Traits and Individual Differences
\item Cognitive Function Executive Processes and Brain Health
\item Motivational Processes Goal Pursuit and Engagement
\item Healthcare Seeking Treatment Adherence and Engagement
\item Emotional Regulation and Pain Affect Processing
\item Mental Health Comorbidity and Psychological Wellbeing
\end{itemize}

\subsection*{Social Subdomains}
\begin{itemize}
\item Social Support Network and Interpersonal Resources
\item Work Disability Occupational Function and Productivity
\item Family Caregiver and Household Dynamics
\item Socioeconomic Status and Health Inequity
\item Healthcare Access Navigation and System Factors
\item Cultural Ethnic and Demographic Context
\item Community Participation and Social Role Functioning
\item Legal Compensation and Medicolegal Systems
\item Health Literacy Education and Patient Empowerment
\item Stigma Social Isolation and Exclusion
\item Return to Work Vocational Rehabilitation and Employment
\item Social Determinants of Pain and Environment
\end{itemize}

\section*{Outputs Produced from This Scheme}
\begin{itemize}
\item Record-level domain loadings.
\item Record-level subdomain loadings.
\item Domain summary and subdomain summary tables.
\item Pairwise loading analyses.
\item Dominance-by-review-type summaries.
\end{itemize}

\section*{Interpretive Note}
This dossier is included because the project uses ontology prompts as a standardized semantic coding scaffold. It should be understood as an analytic coding scheme rather than a human extraction codebook.

\end{document}
